\subsection{Generator Delegation}

Generator delegation means that one generator temporarily hands part of its iteration work to another iterable or generator. Instead of writing an explicit loop that repeatedly yields values from another source, Python provides \texttt{yield from}.

At the simple level, \texttt{yield from} means:

\begin{verbatim}
for item in source:
    yield item
\end{verbatim}

At the deeper level, when the delegated source is itself a generator, \texttt{yield from} also propagates completion values, sent values, thrown exceptions, and close requests through the generator chain.

This subsection starts with the ordinary iteration behavior and then moves to completion and exception propagation.

\subsubsection{\texttt{yield from} as Delegated Iteration over Subgenerators and Iterables}

Start with a generator that manually yields values from a list.

\begin{verbatim}
def manual_quotes():
    quotes = ["D'oh!", "No soup for you!", "The answer is 42."]

    for quote in quotes:
        yield quote

for item in manual_quotes():
    print(item)
\end{verbatim}

Possible output:

\begin{verbatim}
D'oh!
No soup for you!
The answer is 42.
\end{verbatim}

The generator receives values from the list and yields them one by one.

The same behavior can be written with \texttt{yield from}.

\begin{verbatim}
def delegated_quotes():
    quotes = ["D'oh!", "No soup for you!", "The answer is 42."]

    yield from quotes

for item in delegated_quotes():
    print(item)
\end{verbatim}

Possible output:

\begin{verbatim}
D'oh!
No soup for you!
The answer is 42.
\end{verbatim}

The line:

\begin{verbatim}
yield from quotes
\end{verbatim}

means:

\begin{verbatim}
take values from quotes
yield each one outward through this generator
\end{verbatim}

The delegated object can be any iterable, not only another generator.

A list works.

\begin{verbatim}
def from_list():
    yield from ["Homer", "Bart", "Lisa"]

for item in from_list():
    print(item)
\end{verbatim}

Possible output:

\begin{verbatim}
Homer
Bart
Lisa
\end{verbatim}

A tuple works.

\begin{verbatim}
def from_tuple():
    yield from ("Arthur", "Patsy", "Galahad")

for item in from_tuple():
    print(item)
\end{verbatim}

Possible output:

\begin{verbatim}
Arthur
Patsy
Galahad
\end{verbatim}

A string works because strings are iterable over characters.

\begin{verbatim}
def from_string():
    yield from "Ni!"

for item in from_string():
    print(repr(item))
\end{verbatim}

Possible output:

\begin{verbatim}
'N'
'i'
'!'
\end{verbatim}

A \texttt{range} works too.

\begin{verbatim}
def from_range():
    yield from range(40, 43)

for item in from_range():
    print(item)
\end{verbatim}

Possible output:

\begin{verbatim}
40
41
42
\end{verbatim}

The delegated object must be iterable.

\begin{verbatim}
def bad_delegate():
    yield from 42

gen = bad_delegate()

try:
    print(next(gen))
except TypeError as err:
    print(f"error type: {type(err)}")
    print(f"message:    {err}")
\end{verbatim}

Possible output:

\begin{verbatim}
error type: <class 'TypeError'>
message:    'int' object is not iterable
\end{verbatim}

The integer object is not iterable, so \texttt{yield from} cannot pull values from it.

The iterable object can be inspected before delegation.

\begin{verbatim}
items = ["D'oh!", "No soup for you!"]

print(f"type(items): {type(items)}")
print()

selected_names = [
    "__iter__",
    "__len__",
    "__getitem__",
    "append",
    "pop",
]

for name in selected_names:
    print(f"{name:<14} {name in dir(items)}")
\end{verbatim}

Possible output:

\begin{verbatim}
type(items): <class 'list'>

__iter__       True
__len__        True
__getitem__    True
append         True
pop            True
\end{verbatim}

The important method for \texttt{yield from} is \texttt{\_\_iter\_\_}. The object must be able to provide an iterator.

A generator can also delegate to another generator.

\begin{verbatim}
def small_quotes():
    yield "D'oh!"
    yield "No soup for you!"

def all_quotes():
    yield "start"
    yield from small_quotes()
    yield "end"

for item in all_quotes():
    print(item)
\end{verbatim}

Possible output:

\begin{verbatim}
start
D'oh!
No soup for you!
end
\end{verbatim}

The outer generator \texttt{all\_quotes()} yields \texttt{"start"}, then delegates to \texttt{small\_quotes()}, then resumes after the delegated generator is exhausted.

The structure is:

\begin{verbatim}
all_quotes:
    yield "start"

    delegate to small_quotes:
        yield "D'oh!"
        yield "No soup for you!"

    yield "end"
\end{verbatim}

This can be nested further.

\begin{verbatim}
def simpsons():
    yield "D'oh!"
    yield "Mmm... donuts."

def seinfeld():
    yield "No soup for you!"
    yield "These pretzels are making me thirsty."

def monty_python():
    yield "Ni!"
    yield "It is only a model."

def all_lines():
    yield from simpsons()
    yield from seinfeld()
    yield from monty_python()

for line in all_lines():
    print(line)
\end{verbatim}

Possible output:

\begin{verbatim}
D'oh!
Mmm... donuts.
No soup for you!
These pretzels are making me thirsty.
Ni!
It is only a model.
\end{verbatim}

The outer generator does not need three explicit loops. Each \texttt{yield from} delegates one segment of the stream.

The manual equivalent is:

\begin{verbatim}
def all_lines_manual():
    for line in simpsons():
        yield line

    for line in seinfeld():
        yield line

    for line in monty_python():
        yield line
\end{verbatim}

For plain iteration, this is close to what \texttt{yield from} does. But \texttt{yield from} is deeper than loop shorthand when the delegated object is a generator.

We can inspect the outer generator while it is delegating.

\begin{verbatim}
def child():
    yield "D'oh!"
    yield "No soup for you!"

def parent():
    yield "before"
    yield from child()
    yield "after"

gen = parent()

print(f"created: gen.gi_yieldfrom = {gen.gi_yieldfrom}")

print(next(gen))
print(f"after first yield: gen.gi_yieldfrom = {gen.gi_yieldfrom}")

print(next(gen))
print(f"during delegation: gen.gi_yieldfrom = {gen.gi_yieldfrom}")

print(next(gen))
print(f"still during delegation: gen.gi_yieldfrom = {gen.gi_yieldfrom}")

print(next(gen))
print(f"after delegation: gen.gi_yieldfrom = {gen.gi_yieldfrom}")
\end{verbatim}

Possible output:

\begin{verbatim}
created: gen.gi_yieldfrom = None
before
after first yield: gen.gi_yieldfrom = None
D'oh!
during delegation: gen.gi_yieldfrom = <generator object child at 0x...>
No soup for you!
still during delegation: gen.gi_yieldfrom = <generator object child at 0x...>
after
after delegation: gen.gi_yieldfrom = None
\end{verbatim}

The attribute \texttt{gi\_yieldfrom} shows the current delegated iterator while delegation is active.

The generator object has that attribute.

\begin{verbatim}
def child():
    yield "D'oh!"

def parent():
    yield from child()

gen = parent()

print(f"type(gen): {type(gen)}")
print()

selected_names = [
    "__iter__",
    "__next__",
    "send",
    "throw",
    "close",
    "gi_frame",
    "gi_code",
    "gi_yieldfrom",
]

for name in selected_names:
    print(f"{name:<14} {name in dir(gen)}")
\end{verbatim}

Possible output:

\begin{verbatim}
type(gen): <class 'generator'>

__iter__       True
__next__       True
send           True
throw          True
close          True
gi_frame       True
gi_code        True
gi_yieldfrom   True
\end{verbatim}

The core delegation rule is:

\begin{quote}
\texttt{yield from iterable} exposes the delegated iterable's yielded values as if they were yielded by the outer generator.
\end{quote}

\subsubsection{Propagating Values, Exceptions, and Completion Through Delegated Generator Chains}

The simple loop view of \texttt{yield from} explains ordinary yielded values, but it does not explain completion values.

A generator can return a value. That value is carried by \texttt{StopIteration.value}.

\begin{verbatim}
def child():
    yield "child value"
    return "child finished with 42"

gen = child()

print(next(gen))

try:
    next(gen)
except StopIteration as err:
    print(f"completion value: {err.value!r}")
\end{verbatim}

Possible output:

\begin{verbatim}
child value
completion value: 'child finished with 42'
\end{verbatim}

When \texttt{yield from} delegates to this generator, the outer generator can receive that completion value.

\begin{verbatim}
def child():
    yield "child value"
    return "child finished with 42"

def parent():
    result = yield from child()
    print(f"parent received completion: {result!r}")
    yield "parent after child"

for item in parent():
    print(f"caller received: {item!r}")
\end{verbatim}

Possible output:

\begin{verbatim}
caller received: 'child value'
parent received completion: 'child finished with 42'
caller received: 'parent after child'
\end{verbatim}

The line:

\begin{verbatim}
result = yield from child()
\end{verbatim}

does two things:

\begin{enumerate}
    \item yields all normal values produced by \texttt{child()};
    \item after \texttt{child()} completes, binds its return value to \texttt{result}.
\end{enumerate}

So the child's \texttt{return} value does not go directly to the outside \texttt{for} loop. It becomes the value of the \texttt{yield from} expression inside the parent generator.

A plain iterable such as a list has no generator return value, so the \texttt{yield from} expression evaluates to \texttt{None} after the list is exhausted.

\begin{verbatim}
def parent():
    result = yield from ["D'oh!", "No soup for you!"]
    print(f"result after list delegation: {result!r}")
    yield "done"

for item in parent():
    print(f"caller received: {item!r}")
\end{verbatim}

Possible output:

\begin{verbatim}
caller received: "D'oh!"
caller received: 'No soup for you!'
result after list delegation: None
caller received: 'done'
\end{verbatim}

The list contributes yielded items, but no completion value.

Delegation can be chained.

\begin{verbatim}
def grandchild():
    yield "grandchild value"
    return "grandchild done"

def child():
    result = yield from grandchild()
    print(f"child received: {result!r}")
    return "child done"

def parent():
    result = yield from child()
    print(f"parent received: {result!r}")
    return "parent done"

gen = parent()

while True:
    try:
        item = next(gen)
        print(f"caller received: {item!r}")
    except StopIteration as err:
        print(f"caller sees completion: {err.value!r}")
        break
\end{verbatim}

Possible output:

\begin{verbatim}
caller received: 'grandchild value'
child received: 'grandchild done'
parent received: 'child done'
caller sees completion: 'parent done'
\end{verbatim}

The stream value flows outward:

\begin{verbatim}
grandchild yield -> child -> parent -> caller
\end{verbatim}

The completion values flow one level upward at a time:

\begin{verbatim}
grandchild return -> child receives as yield from result
child return      -> parent receives as yield from result
parent return     -> caller sees as StopIteration.value
\end{verbatim}

This is the main reason \texttt{yield from} is not only loop shorthand.

Exceptions can also move through a delegated chain. If the caller throws an exception into the outer generator while it is delegating, Python sends that exception into the currently delegated generator.

\begin{verbatim}
def child():
    try:
        yield "child ready"
        yield "child still running"
    except ValueError as err:
        print(f"child caught: {err}")
        yield "child recovered"

    return "child done"

def parent():
    result = yield from child()
    print(f"parent received completion: {result}")
    yield "parent done"

gen = parent()

print(next(gen))
print(gen.throw(ValueError("bad value")))
print(next(gen))
\end{verbatim}

Possible output:

\begin{verbatim}
child ready
child caught: bad value
child recovered
parent received completion: child done
parent done
\end{verbatim}

The call:

\begin{verbatim}
gen.throw(ValueError("bad value"))
\end{verbatim}

throws the exception into \texttt{parent()}, but \texttt{parent()} is currently suspended at:

\begin{verbatim}
yield from child()
\end{verbatim}

So the exception is forwarded into \texttt{child()}. The child generator catches it, yields a recovery value, then finishes normally.

The exception object is ordinary.

\begin{verbatim}
err = ValueError("bad value")

print(f"err:       {err}")
print(f"type(err): {type(err)}")
print()

selected_names = [
    "args",
    "__class__",
    "__str__",
    "__repr__",
    "with_traceback",
]

for name in selected_names:
    print(f"{name:<18} {name in dir(err)}")

print()
print(f"err.args: {err.args}")
\end{verbatim}

Possible output:

\begin{verbatim}
err:       bad value
type(err): <class 'ValueError'>

args               True
__class__          True
__str__            True
__repr__           True
with_traceback     True

err.args: ('bad value',)
\end{verbatim}

If the delegated generator does not handle the thrown exception, the exception propagates outward.

\begin{verbatim}
def child():
    yield "child ready"
    yield "child still running"

def parent():
    yield from child()
    yield "parent done"

gen = parent()

print(next(gen))

try:
    gen.throw(ValueError("bad value"))
except ValueError as err:
    print(f"caller caught: {err}")
\end{verbatim}

Possible output:

\begin{verbatim}
child ready
caller caught: bad value
\end{verbatim}

The exception entered the delegated generator, was not caught there, was not caught by the parent generator, and finally reached the caller.

A parent can catch exceptions around the delegation.

\begin{verbatim}
def child():
    yield "child ready"
    yield "child still running"

def parent():
    try:
        yield from child()
    except ValueError as err:
        print(f"parent caught: {err}")
        yield "parent recovered"

    yield "parent done"

gen = parent()

print(next(gen))
print(gen.throw(ValueError("bad value")))
print(next(gen))
\end{verbatim}

Possible output:

\begin{verbatim}
child ready
parent caught: bad value
parent recovered
parent done
\end{verbatim}

Here the child did not handle the exception. The parent handled it after the delegated chain failed.

The propagation structure is:

\begin{verbatim}
caller throws into parent
    parent is delegating to child
        exception is offered to child
        if child handles it:
            child continues
        else:
            exception returns to parent
            parent may handle it
        else:
            exception reaches caller
\end{verbatim}

Completion and exceptions are therefore routed through the delegation chain in a structured way.

A final example combines yielded values, completion values, and exception handling.

\begin{verbatim}
def child():
    try:
        yield "child: first"
        yield "child: second"
    except RuntimeError as err:
        print(f"child caught: {err}")
        yield "child: recovered"

    return "child completion value"

def parent():
    print("parent before delegation")

    result = yield from child()

    print(f"parent got completion: {result!r}")
    yield "parent: after child"

    return "parent completion value"

gen = parent()

print(next(gen))
print(next(gen))
print(gen.throw(RuntimeError("something happened")))

try:
    print(next(gen))
    print(next(gen))
except StopIteration as err:
    print(f"caller got completion: {err.value!r}")
\end{verbatim}

Possible output:

\begin{verbatim}
parent before delegation
child: first
child: second
child caught: something happened
child: recovered
parent got completion: 'child completion value'
parent: after child
caller got completion: 'parent completion value'
\end{verbatim}

The important events are:

\begin{enumerate}
    \item \texttt{parent()} starts and delegates to \texttt{child()};
    \item \texttt{child()} yields values outward through \texttt{parent()};
    \item the caller throws an exception into the chain;
    \item \texttt{child()} catches the exception and yields a recovery value;
    \item \texttt{child()} returns a completion value;
    \item \texttt{parent()} receives that value as the result of \texttt{yield from};
    \item \texttt{parent()} yields one more value;
    \item \texttt{parent()} returns its own completion value;
    \item the caller sees the final completion value through \texttt{StopIteration.value}.
\end{enumerate}

The compact rule is:

\begin{quote}
\texttt{yield from} delegates a stream of yielded values outward and routes completion values, thrown exceptions, and generator control operations through the delegated iterator. It is a generator-chain control mechanism, not merely a shorter \texttt{for} loop.
\end{quote}